% Seccion retirada del Apendice C (appendices/bsplines.tex) el 2026-10-13, a pedido del autor:
% es un detalle de implementacion cuya explicacion completa no cabe en el manuscrito. Se conserva
% aqui para el documento de detalles adicionales. Los numeros salen de docs/tablas/valores_texto.md
% (etapa 3, tesis/diagnostico_estados_espurios.jl). La limitacion sigue documentada en la seccion
% de limitaciones de chapters/resultados.tex (\label{sec:espacio_ci}).

\section{El estado espurio del espacio de orbitales del CI}
\label{sec:estado_espurio}

El espacio de orbitales del cálculo de interacción de configuraciones se obtiene diagonalizando, para cada $l$, el operador de Fock de core congelado, y a diferencia del ciclo autoconsistente esa diagonalización emplea todos los splines, incluido el primero, que no se anula en el origen. En los canales con $l > 0$ la barrera centrífuga penaliza a ese spline y no ocurre nada. En el canal $s$, en cambio, una función que no se anula en $r = 0$ tiene una energía potencial coulombiana no acotada inferiormente, y la cuadratura la regulariza en un autovalor espurio muy por debajo de cualquier estado físico: $-1178.6$ Hartree en el carbono y $-108515.4$ en el estaño, frente a una cota hidrogenoide de $-Z^2/2$. El motor descarta cualquier autovalor menor que $-Z^2$, y el filtro elimina exactamente ese estado en los cuatro elementos.

El filtro no repara el resto del espectro. Los demás autovectores del canal $s$ son ortogonales al estado espurio y no a los orbitales $s$ del core que produjo el ciclo autoconsistente, y los virtuales $s$ del espacio activo conservan un solapamiento con esos orbitales de hasta $3.6\times10^{-2}$, con un peso máximo del core de $2.0\times10^{-3}$ en un solo virtual. En el canal $p$ el solapamiento no pasa de $8.8\times10^{-4}$, y en el canal $d$ del germanio y del estaño, donde solo refleja la diferencia entre el operador de core congelado y el del ciclo, llega a $8.5\times10^{-3}$. El cálculo de interacción de configuraciones supone que el espacio activo es ortogonal al core congelado, de modo que esta es una aproximación de la implementación. La desviación del canal $s$, la mayor, afecta solo a las configuraciones con electrones $s$, que entran en los términos $^1S$ y $^1D$ y no en el $^3P$, y su efecto sobre los singletes no se cuantificó (Sección \ref{sec:limitaciones}).
